% Phụ lục - Bảng phân công lấy từ mục 3 của README
\section{BẢNG PHÂN CÔNG VÀ ĐÓNG GÓP CỦA THÀNH VIÊN}\label{app:phan-cong}

\begin{table}[H]
\centering
\caption{Phân công nhiệm vụ và đóng góp của các thành viên}
\label{tab:phan-cong}
\small
\begin{tabularx}{\textwidth}{|>{\raggedright\arraybackslash}p{2.4cm}|>{\raggedright\arraybackslash}X|*{4}{>{\centering\arraybackslash}p{1.45cm}|}}
\hline
\rowcolor{LightCyan}
\textbf{Thành viên} & \textbf{Nội dung phụ trách} & \textbf{Báo cáo (trang)} & \textbf{Slide} & \textbf{Thuyết trình (phút)} & \textbf{Đóng góp} \\
\hline
Trần Bá Đạt & Chương~\ref{chap:mo-dau} và Chương~\ref{chap:co-so-ly-thuyet}: mở đầu, cơ sở lý thuyết, khảo sát công trình liên quan, bảng so sánh; biên tập chung; dẫn dắt buổi thuyết trình & 8--9 & $\sim$8 & 4--5 & \dots\% \\
\hline
Vũ Đức Minh & Chương~\ref{chap:danh-gia}, Chương~\ref{chap:chuyen-doi-so} và Chương~\ref{chap:ket-luan}: đánh giá tổng hợp, tác động chuyển đổi số và thách thức, kết luận & 5--7 & $\sim$6 & 4--5 & \dots\% \\
\hline
Phạm Tuấn Anh & Chương~\ref{chap:thuc-nghiem}: dữ liệu, độ chính xác, khả năng tổng quát hoá, tốc độ, thử nghiệm thực tế & 5--6 & $\sim$6 & 4--5 & \dots\% \\
\hline
Nguyễn Đồng Hoàng & Chương~\ref{chap:fanci}: kiến trúc 3 module, 21 đặc trưng, RF/SVM & 6--7 & $\sim$7 & 4--5 & \dots\% \\
\hline
\end{tabularx}
\end{table}

\canviet{Bổ sung vai trò trong thực nghiệm minh hoạ (nếu có) và tỷ lệ đóng góp đã được cả nhóm thống nhất.}

\clearpage
\section{MÃ NGUỒN THỰC NGHIỆM MINH HOẠ}\label{app:ma-nguon}
\canviet{Tuỳ chọn. Mã nguồn script Python ở Mục~\ref{sec:thuc-nghiem-nhom}, hoặc đường dẫn tới repository GitHub của nhóm.}
